% ID: SEL-D0020-OB-E02
% ESTADO: validada. S021, revisada con los criterios oficiales.
\textbf{a)} En $-1$, el valor y límite izquierdo son $1-a$,
mientras que el límite derecho vale $-1$. La continuidad exige
$a=2$. En $0$ ambos límites y el valor son cero para todo $b$.

Una vez asegurada la continuidad, igualamos las pendientes laterales.
Las derivadas de los tres trozos son $a$, $2/(x+2)^2$ y $2x-b$.
En $-1$ las pendientes son $a$ y $2$, que coinciden para $a=2$.
En $0$ son $1/2$ y $-b$, que coinciden para $b=-1/2$.
Por tanto $\boxed{a=2,\ b=-1/2}$.

\textbf{b)} Con estos valores, las derivadas son $2$,
$2/(x+2)^2$ y $2x+1/2$ en los respectivos interiores.
Son positivas en cada tramo y en ambos empalmes. Por el signo de
la derivada, $f$ es $\boxed{\text{estrictamente creciente en }\mathbb R}$.
